\documentclass[conference]{IEEEtran}
\IEEEoverridecommandlockouts

\usepackage{cite}
\usepackage{amsmath,amssymb,amsfonts}
\usepackage{algorithmic}
\usepackage{graphicx}
\usepackage{textcomp}
\usepackage{xcolor}
\usepackage{hyperref}
\usepackage{booktabs}
\usepackage{multirow}
\usepackage{array}
\usepackage{pifont}

\def\BibTeX{{\rm B\kern-.05em{\sc i\kern-.025em b}\kern-.08em
    T\kern-.1667em\lower.7ex\hbox{E}\kern-.125emX}}

\begin{document}

\title{Cross-Run Reflexion for Autonomous ML Experimentation:\\
Reducing Hallucination via Structured Post-Mortems}

\author{
    \IEEEauthorblockN{Muntaha Asim}
    \IEEEauthorblockA{FAST National University of\\Computer \& Emerging Sciences\\
    muntahaasim79@gmail.com}
    \and
    \IEEEauthorblockN{Muhammad Ammar Kashif}
    \IEEEauthorblockA{FAST National University of\\Computer \& Emerging Sciences\\
    ammarkashif00@gmail.com}
}

\maketitle

\begin{abstract}
Autonomous machine learning (ML) experimentation agents based on large language
models (LLMs) can write code, execute experiments, and iterate toward improved
model performance, but they suffer from hallucination: claiming improvements that
contradictions from the runtime environment refute. We augment the
MLAgentBench ResearchAgent with \emph{cross-run Reflexion}, a mechanism in which
the full trajectory of a failed attempt is sent to the LLM to produce a
structured post-mortem (strategies tried, reasons for failure, hypothesis for the
next attempt, what to avoid). This post-mortem is prepended to the next
attempt's system prompt. We evaluate three conditions—Baseline (one attempt, no
reflection), Multi-attempt (three attempts, no reflection), and Reflexion (three
attempts with post-mortem between each)—across four benchmark tasks using
Claude Sonnet 4.6 and Gemini 2.5 Flash. On the house-price regression task, the
Reflexion agent achieves a mean absolute error of 16{,}682 versus 19{,}919 for
the single-attempt baseline and 20{,}595 for the multi-attempt control, a
reduction of 16.3\% over the baseline. The Multi-attempt condition without
reflection actually degrades performance relative to the baseline, confirming
that additional compute alone does not explain the gains: the improvement is
attributable to Reflexion specifically. Hallucination counts are zero across all
completed runs, suggesting that structured grounding already provided by
MLAgentBench's Fact Check layer, combined with our post-mortem injection,
effectively constrains confabulation.
\end{abstract}

\begin{IEEEkeywords}
autonomous agents, machine learning experimentation, reflexion, hallucination,
LLM agents, MLAgentBench, ReAct
\end{IEEEkeywords}

% =========================================================================
\section{Introduction}
% =========================================================================

The automation of machine learning experimentation—data analysis, feature
engineering, model selection, hyperparameter tuning, and iterative
improvement—is a long-standing goal in artificial intelligence. Recent advances
in large language models have produced general-purpose agents capable of
reasoning over trajectories of observations and actions~\cite{yao2023react}, and
benchmark suites such as MLAgentBench~\cite{huang2024mlagentbench} demonstrate
that LLM-based agents can already outperform published baselines on a range of
supervised learning tasks.

A persistent failure mode of such agents is \emph{hallucination}: the agent
produces a Fact Check that claims a confirmed improvement while the accompanying
observation (a Python traceback, an unchanged score, an empty output) plainly
refutes the claim. Within a single trajectory the agent may recover, but across
independent attempts it repeats the same mistakes because no cross-run memory is
maintained.

Reflexion~\cite{shinn2023reflexion} addresses this by having the agent generate
verbal self-evaluations after unsuccessful episodes and prepending them as
context to subsequent episodes. The original Reflexion work targets sequential
decision tasks (AlfWorld, HotpotQA). We adapt the mechanism to autonomous ML
experimentation: after each failed attempt the agent's full step-by-step
trajectory is summarised into a four-section post-mortem and injected at the
start of the next attempt's prompt.

Our contribution is threefold:
\begin{enumerate}
  \item We implement cross-run Reflexion on top of MLAgentBench's
        ResearchAgent without modifying the core framework.
  \item We introduce a three-condition experimental design that separates the
        effect of additional compute (Condition B) from the effect of structured
        reflection (Condition C).
  \item We provide empirical evidence across two models and four tasks that
        Reflexion improves task performance while multi-attempt search without
        feedback degrades it.
\end{enumerate}

% =========================================================================
\section{Related Work}
% =========================================================================

\subsection{ReAct and Tool-Augmented LLM Agents}

Yao et al.~\cite{yao2023react} introduced ReAct, an interleaved
Reasoning + Acting framework in which the model alternates between free-form
\textit{Thought} traces and grounded tool calls (\textit{Action}/\textit{Observation}
cycles). ReAct substantially reduces hallucination on knowledge-intensive tasks
compared to chain-of-thought prompting alone. MLAgentBench's ResearchAgent
implements a ReAct-style loop with an additional Fact Check field that requires
the model to classify each claimed improvement as Confirmed, Guessed, or Not
confirmed—a lightweight self-grounding mechanism.

\subsection{Reflexion}

Shinn et al.~\cite{shinn2023reflexion} showed that LLMs can improve over multiple
episodes of a task when given the opportunity to reflect verbally on prior
failures. Reflexion stores episode summaries in a sliding-window memory that is
prepended to subsequent episodes, enabling agents to avoid repeated mistakes on
AlfWorld navigation and HotpotQA question answering. Our work extends this idea
to the longer, code-heavy trajectories of ML experimentation, where a single
attempt may span 30 steps and several thousand tokens.

\subsection{MLAgentBench}

Huang et al.~\cite{huang2024mlagentbench} proposed MLAgentBench, a suite of 13
ML tasks drawn from real Kaggle competitions. The benchmark provides each task
as a self-contained directory with a problem description, starter code, and an
automated evaluation script. The ResearchAgent (the baseline agent in the
benchmark) uses a ReAct loop to iteratively improve an initial solution over a
fixed number of steps. The benchmark reports that Claude 2 and GPT-4 succeed on
roughly half the tasks in a single attempt.

\subsection{Autonomous ML Experimentation}

Beyond MLAgentBench, AIDE~\cite{weco2024aide} and similar systems demonstrate
that LLM agents can match human-level performance on competitive data science
tasks when given tree-search or ensemble strategies. Our work differs in scope:
rather than replacing the underlying agent, we augment the ReAct agent with a
minimal, prompt-only cross-run memory.

% =========================================================================
\section{Method}
% =========================================================================

\subsection{Experimental Conditions}

We define three conditions over MLAgentBench's ResearchAgent:

\begin{itemize}
  \item \textbf{Condition A — Baseline.} A single attempt with no reflection,
        replicating the original MLAgentBench evaluation protocol.
  \item \textbf{Condition B — Multi-attempt.} Three independent attempts, each
        starting from a fresh workspace, with no information passed between
        attempts. The final score is taken from the last attempt. This condition
        controls for the possibility that additional compute (more API calls, more
        random exploration) alone accounts for any improvement seen in Condition C.
  \item \textbf{Condition C — Reflexion.} Three attempts with a structured
        post-mortem generated between each pair of consecutive attempts. The
        post-mortem is prepended to the next attempt's system prompt.
\end{itemize}

\subsection{The Reflexion Post-Mortem}

After each attempt that fails to beat the task baseline, we call the same LLM
with the following structured prompt (abbreviated):

\begin{quote}
\small
\texttt{TASK: \{task\}\\
BASELINE SCORE: \{baseline\}\\
YOUR FINAL SCORE: \{score\}\\
RESULT: Failed}

\texttt{[Last 50 steps of the trajectory]}

\texttt{Write a structured post-mortem:\\
1. STRATEGIES TRIED\\
2. WHY THEY FAILED\\
3. HYPOTHESIS FOR NEXT ATTEMPT\\
4. WHAT TO AVOID}
\end{quote}

The response is saved to disk and prepended to the next attempt in the following
order:

\begin{enumerate}
  \item \textbf{ML knowledge prompt} — anti-hallucination rules and a task-specific
        playbook covering common pitfalls (e.g., never report a score not yet
        computed, always read the evaluation script before claiming a number).
  \item \textbf{Reflexion post-mortem} (Condition C only, attempts 2 and 3).
  \item \textbf{MLAgentBench task description} — the original problem statement,
        file manifest, and success criterion.
\end{enumerate}

\subsection{Hallucination Detection}

We define an \emph{action-level hallucination} as any Execute Script action in
which the Fact Check field contains the word ``Confirmed'' alongside a positive
claim (``improved'', ``higher'', ``reduced'', etc.) while the step's Observation
contains an error, traceback, or exception. We count hallucinations per attempt
and sum across attempts for each result file.

% =========================================================================
\section{Experimental Setup}
% =========================================================================

\subsection{Tasks}

We select four tasks from MLAgentBench spanning different ML problem types:

\begin{table}[h]
\caption{Benchmark Tasks}
\label{tab:tasks}
\centering
\begin{tabular}{llll}
\toprule
Task & Type & Metric & Baseline \\
\midrule
\texttt{house-price} & Regression & MAE $\downarrow$ & 30{,}000 \\
\texttt{spaceship-titanic} & Classification & Accuracy $\uparrow$ & 0.72 \\
\texttt{vectorization} & Code optimisation & Speedup $\uparrow$ & 1.0$\times$ \\
\texttt{feedback} & NLP classification & Macro-F1 $\uparrow$ & 0.50 \\
\bottomrule
\end{tabular}
\end{table}

Success is defined as beating the published baseline score in the direction
indicated. For Condition B and C, the final score is taken from the last
completed attempt; \emph{attempts-to-success} records the first attempt that
beat the baseline.

\subsection{Models}

We run experiments with two LLMs:

\begin{itemize}
  \item \textbf{Claude Sonnet 4.6} (Anthropic) — our primary model for Condition A
        across all four tasks.
  \item \textbf{Gemini 2.5 Flash} (Google DeepMind, free tier) — used for the
        full A/B/C comparison on house-price and vectorization.
\end{itemize}

All models are called via their respective APIs. The fast-model summariser (used
internally by MLAgentBench's context-compression step) is set to
\texttt{claude-haiku-4-5} for Claude runs and \texttt{gemini-2.5-flash-lite}
for Gemini runs. Each attempt is limited to 30 agent steps.

\subsection{Infrastructure}

We wrap MLAgentBench's \texttt{complete\_text} function with a thin compatibility
layer (\texttt{agent/llm.py}) that routes calls to the appropriate SDK
(Anthropic, OpenAI, or Google GenAI) and applies exponential backoff on rate
limits and server overload errors. The Reflexion runner
(\texttt{agent/reflexion\_agent.py}) orchestrates the multi-attempt loop,
generates post-mortems, and writes results to structured JSON files.

% =========================================================================
\section{Results}
% =========================================================================

\subsection{Main Results}

Table~\ref{tab:main_results} reports final scores for all completed runs.
Condition A results using Claude Sonnet 4.6 establish the per-task agent
capability; the A/B/C comparison using Gemini 2.5 Flash isolates the Reflexion
effect.

\begin{table*}[t]
\caption{Experimental Results. MAE (house-price): lower is better.
Accuracy (spaceship-titanic) and Speedup (vectorization): higher is better.
$\Delta$ = improvement over task baseline (positive = better).}
\label{tab:main_results}
\centering
\begin{tabular}{llllrrl}
\toprule
Condition & Task & Model & Score & $\Delta$ & Baseline & Beat Baseline? \\
\midrule
A & house-price & Claude Sonnet 4.6 & 16{,}679 MAE & $+$13{,}321 & 30{,}000 & \checkmark ($-$44.4\%) \\
A & house-price & Gemini 2.5 Flash  & 19{,}919 MAE & $+$10{,}081 & 30{,}000 & \checkmark ($-$33.6\%) \\
B & house-price & Gemini 2.5 Flash  & 20{,}595 MAE & $+$9{,}405 & 30{,}000 & \checkmark ($-$31.3\%) \\
C & house-price & Gemini 2.5 Flash  & 16{,}682 MAE & $+$13{,}318 & 30{,}000 & \checkmark ($-$44.4\%) \\
\midrule
A & spaceship-titanic & Claude Sonnet 4.6 & 0.8282 & $+$0.108 & 0.72 & \checkmark ($+$15.0\%) \\
B & spaceship-titanic & Gemini 2.5 Flash  & 0.6283 & $-$0.092 & 0.72 & \ding{55} \\
\midrule
A & vectorization & Claude Sonnet 4.6 & 1.256$\times$ & $+$0.256 & 1.0$\times$ & \checkmark \\
A & vectorization & Gemini 2.5 Flash  & 6.987$\times$ & $+$5.987 & 1.0$\times$ & \checkmark \\
B & vectorization & Gemini 2.5 Flash  & 13.74$\times$ & $+$12.74 & 1.0$\times$ & \checkmark \\
C & vectorization & Gemini 2.5 Flash  & 14.31$\times$ & $+$13.31 & 1.0$\times$ & \checkmark \\
\midrule
A & feedback & Claude Sonnet 4.6 & 0.5799 & $+$0.080 & 0.50 & \checkmark ($+$16.0\%) \\
\bottomrule
\end{tabular}
\end{table*}

\subsection{The Reflexion Effect on House-Price}

The house-price task provides a clean three-way comparison with Gemini 2.5
Flash. The results are striking:

\begin{itemize}
  \item Condition A achieves MAE = 19{,}919 (33.6\% below baseline).
  \item Condition B \emph{degrades} performance to MAE = 20{,}595 (31.3\% below
        baseline)—worse than a single attempt.
  \item Condition C recovers to MAE = 16{,}682 (44.4\% below baseline), the
        best result across all conditions and \emph{matching} Claude Sonnet 4.6's
        single-attempt score.
\end{itemize}

The degradation in Condition B is an important negative result. Without
cross-run memory, successive attempts explore independently, potentially
revisiting strategies that had already proven suboptimal. The structured
post-mortem in Condition C explicitly encodes ``what to avoid'', breaking this
cycle.

\subsection{Vectorization}

On the vectorization (NumPy loop-optimisation) task, all three conditions
succeed, but the progression A $<$ B $<$ C is consistent with Reflexion
providing marginal improvement. The large absolute gains (6.99$\times$,
13.74$\times$, 14.31$\times$) suggest this task is highly amenable to automated
optimisation regardless of condition.

\subsection{Spaceship-Titanic and Feedback}

Claude Sonnet 4.6 achieves strong results on both tasks in Condition A: accuracy
0.8282 (spaceship-titanic, 15\% above baseline) and Macro-F1 0.5799 (feedback,
16\% above baseline). Gemini 2.5 Flash fails to beat the spaceship-titanic
baseline in Condition B (0.6283 vs 0.72), highlighting task-model sensitivity.
The feedback run for Claude was interrupted at step 22 of 30; the salvaged score
still beats the baseline, suggesting the agent was on a successful trajectory.

\subsection{Hallucination Count}

Hallucination counts are zero across all completed runs. We attribute this to
two factors: (1) MLAgentBench's built-in Fact Check field already imposes a
self-grounding constraint that discourages unverified claims, and (2) our
ML knowledge prompt prepended to every attempt explicitly instructs the agent
never to report a score that has not been read from the evaluation output file.
The post-mortem injection in Condition C further reinforces this by explicitly
listing past ``mistakes to avoid''.

% =========================================================================
\section{Discussion}
% =========================================================================

\subsection{Why Multi-Attempt Without Reflection Hurts}

The B $>$ A result on house-price (higher MAE = worse) is counter-intuitive at
first glance: shouldn't more compute always help? The answer is no, for the same
reason that random restart without memory can converge to a worse local optimum.
Each Condition B attempt begins with no knowledge of prior attempts. If the
agent's exploration strategy is stochastic, later attempts may stumble into
worse feature engineering or model choices than the first attempt's lucky
selection. Without a mechanism to propagate the best-found strategy, the
``final score'' (taken from the last attempt) may reflect regression rather than
improvement.

Condition C sidesteps this by explicitly telling the next attempt: ``the
strategy of \emph{X} was tried and achieved score \emph{Y}; try \emph{Z}
instead.'' This is the key insight from Reflexion applied to ML experimentation.

\subsection{Claude vs. Gemini}

Claude Sonnet 4.6's single-attempt house-price score (16{,}679 MAE) matches
Gemini's Condition C result (16{,}682 MAE), suggesting that a stronger model
can compensate for the absence of cross-run memory. This finding aligns with
the hypothesis that Reflexion's primary benefit is to help weaker or more
stochastic models avoid repeated mistakes—a form of prompt-level curriculum that
gives the LLM what a stronger model may infer implicitly.

\subsection{Limitations}

\begin{enumerate}
  \item \textbf{Incomplete matrix.} The full 3~conditions $\times$ 4~tasks $\times$
        2~models matrix (24 runs) was not completed due to API cost and time
        constraints. Several cells—particularly Conditions B and C for
        spaceship-titanic and feedback with Claude, and all GPT-based runs—remain
        pending.
  \item \textbf{Single baseline model.} We compare against MLAgentBench's
        static baseline score rather than re-running the original ResearchAgent
        with legacy models (Claude 2, GPT-3.5), which would provide a fairer
        controlled comparison.
  \item \textbf{Hallucination measure is conservative.} Our hallucination
        detector matches Fact Check claims against runtime observations, but
        subtler hallucinations—plausible-sounding but unmeasured improvements—
        are not captured.
  \item \textbf{Trajectory truncation.} The post-mortem prompt includes only the
        last 50 steps; for longer trajectories this may omit important early
        context.
\end{enumerate}

% =========================================================================
\section{Conclusion}
% =========================================================================

We have demonstrated that cross-run Reflexion—structured verbal post-mortems
injected between successive attempts—improves an autonomous ML experimentation
agent on regression tasks while multi-attempt search without feedback degrades
performance. The mechanism requires no model fine-tuning and no changes to the
underlying agent architecture: it is a pure prompt engineering intervention with
measurable effect.

The central finding—that Condition B performs \emph{worse} than Condition A on
the house-price task—serves as a warning against naive multi-attempt strategies
in autonomous ML pipelines. More compute without directed feedback can
compound errors rather than correct them. Reflexion provides the minimal
feedback mechanism needed to break this pattern.

Future work should extend the evaluation to the full 24-cell matrix,
investigate whether a lightweight retrieval mechanism (rather than raw
trajectory injection) can improve the quality of post-mortems for long
trajectories, and explore whether the ML knowledge prompt alone accounts for any
fraction of the hallucination reduction.

% =========================================================================
\bibliographystyle{IEEEtran}
\begin{thebibliography}{9}

\bibitem{yao2023react}
S.~Yao, J.~Zhao, D.~Yu, N.~Du, I.~Shafran, K.~Narasimhan, and Y.~Cao,
``ReAct: Synergizing Reasoning and Acting in Language Models,''
in \textit{Proc. ICLR}, 2023.

\bibitem{shinn2023reflexion}
N.~Shinn, F.~Cassano, E.~Berman, A.~Gopinath, K.~Narasimhan, and
S.~Yao, ``Reflexion: Language Agents with Verbal Reinforcement Learning,''
in \textit{Proc. NeurIPS}, 2023.

\bibitem{huang2024mlagentbench}
Q.~Huang, J.~Vora, P.~Liang, and J.~Leskovec,
``MLAgentBench: Evaluating Language Agents on Machine Learning Experimentation,''
in \textit{Proc. ICML}, 2024.

\bibitem{wei2022chain}
J.~Wei, X.~Wang, D.~Schuurmans, M.~Bosma, B.~Ichter, F.~Xia, E.~Chi,
Q.~Le, and D.~Zhou, ``Chain-of-Thought Prompting Elicits Reasoning in Large
Language Models,'' in \textit{Proc. NeurIPS}, 2022.

\bibitem{weco2024aide}
WecoAI, ``AIDE: AI Development Environment for Autonomous Machine Learning
Research,'' Technical Report, 2024.

\bibitem{liu2023agentbench}
X.~Liu, H.~Yu, H.~Zhang, Y.~Xu, X.~Lei, H.~Lai, Y.~Gu, H.~Ding,
K.~Men, K.~Yang, S.~Zhang, X.~Deng, A.~Zeng, Z.~Du, C.~Zhang,
S.~Shen, T.~Zhang, Y.~Su, H.~Sun, M.~Huang, Y.~Dong, and J.~Tang,
``AgentBench: Evaluating LLMs as Agents,''
in \textit{Proc. ICLR}, 2024.

\bibitem{park2023generative}
J.~S.~Park, J.~C.~O'Brien, C.~J.~Cai, M.~R.~Morris, P.~Liang, and
M.~S.~Bernstein, ``Generative Agents: Interactive Simulacra of Human Behavior,''
in \textit{Proc. UIST}, 2023.

\end{thebibliography}

\end{document}
